% =====================================================================
%  plan-de-proyecto.tex  --  Plan de Proyecto de LEIP
%  ARCHIVO UNICO: no necesita ningun otro archivo.
%  Compilar en Overleaf con pdfLaTeX (dos pasadas por el indice).
%  Estilo consistente con casos-de-uso.tex
% =====================================================================
\documentclass[11pt,a4paper]{article}

% ---------------------------------------------------------------------
%  DATOS DEL PROYECTO
% ---------------------------------------------------------------------
\newcommand{\proyectoNombre}{LEIP}
\newcommand{\proyectoNombreLargo}{LatAm Economic Intelligence Platform}
\newcommand{\proyectoUniversidad}{Universidad Cat\'olica del Uruguay}
\newcommand{\proyectoFacultad}{Facultad de Ingenier\'ia y Tecnolog\'ias}
\newcommand{\proyectoCarrera}{Ingenier\'ia en Inform\'atica}
\newcommand{\docVersion}{1.2}
\newcommand{\docTitulo}{Plan de Proyecto}
\newcommand{\docFecha}{Septiembre 2026}

% ---------------------------------------------------------------------
%  PREAMBULO
% ---------------------------------------------------------------------
\usepackage[utf8]{inputenc}
\usepackage[T1]{fontenc}
\IfFileExists{lmodern.sty}{\usepackage{lmodern}\usepackage{microtype}}{}

\IfFileExists{spanish.ldf}{%
  \usepackage[spanish,es-tabla,es-noquoting]{babel}%
}{%
  \usepackage[english]{babel}%
  \renewcommand{\contentsname}{\'Indice}%
  \renewcommand{\listfigurename}{\'Indice de figuras}%
  \renewcommand{\listtablename}{\'Indice de tablas}%
  \renewcommand{\tablename}{Tabla}%
  \renewcommand{\figurename}{Figura}%
}

\usepackage[a4paper,margin=2.5cm,headheight=15pt]{geometry}
\usepackage{fancyhdr}
\usepackage{parskip}
\usepackage{booktabs}
\usepackage{tabularx}
\usepackage{longtable}
\usepackage{array}
\usepackage{enumitem}
\setlist{noitemsep,topsep=4pt}
\usepackage{graphicx}
\usepackage{xcolor}
\usepackage[hidelinks]{hyperref}
\usepackage{url}
\urlstyle{same}

\pagestyle{fancy}
\fancyhf{}
\fancyhead[L]{\small\proyectoNombre{} --- \docTitulo}
\fancyhead[R]{\small v\docVersion}
\fancyfoot[C]{\small\thepage}
\renewcommand{\headrulewidth}{0.4pt}

\usepackage{titlesec}
\titleformat{\section}{\Large\bfseries}{\thesection}{0.6em}{}
\titleformat{\subsection}{\large\bfseries}{\thesubsection}{0.6em}{}
\titleformat{\subsubsection}{\normalsize\bfseries}{\thesubsubsection}{0.6em}{}

% --- Portada -----------------------------------------------------------
\newcommand{\logoUniversidad}{img/logo-ucu.png}
\newcommand{\portada}[1]{%
  \begin{titlepage}
  \centering
  \vspace*{1.5cm}
  \IfFileExists{\logoUniversidad}{%
    \includegraphics[width=0.40\textwidth]{\logoUniversidad}\par
    \vspace{1.2cm}%
  }{%
    \fbox{\parbox[c][3cm][c]{0.5\textwidth}{\centering\small\texttt{\logoUniversidad}}}\par
    \vspace{1.2cm}%
  }
  {\Large \proyectoUniversidad\par}
  {\large \proyectoFacultad\par}
  {\large \proyectoCarrera\par}
  \vspace{3cm}
  {\Huge\bfseries #1\par}
  \vspace{0.8cm}
  {\Large \proyectoNombre{} --- \proyectoNombreLargo\par}
  \vspace{2cm}
  {\large Versi\'on \docVersion\par}
  {\large \docFecha\par}
  \vfill
  \begin{tabular}{c}
    Sebastian Forische \\
    Armando Hernandez de la Vega \\
    Luana Acu\~na \\
    Bruno Sebastian Olivera Viera Gomez \\
  \end{tabular}
  \vspace{1.5cm}
  \end{titlepage}
}

% --- Historial de revisiones -------------------------------------------
\newenvironment{historial}{%
  \section*{Historial de revisiones}
  \addcontentsline{toc}{section}{Historial de revisiones}
  \begin{tabular}{@{}l l p{0.52\textwidth} l@{}}
  \toprule
  \textbf{Versi\'on} & \textbf{Fecha} & \textbf{Descripci\'on} & \textbf{Autor} \\
  \midrule
}{%
  \bottomrule
  \end{tabular}
  \bigskip
}

\begin{document}

\portada{\docTitulo}

\tableofcontents
\bigskip

\begin{historial}
1.0 & 31/08/2026 & Versi\'on inicial del Plan de Proyecto & Equipo LEIP \\
1.1 & 11/09/2026 & Alineaci\'on terminol\'ogica con el Documento de Visi\'on: se retira el t\'ermino MVP y se unifica la referencia a la Lista de Riesgos & Equipo LEIP \\
1.2 & 12/09/2026 & Revisi\'on de redacci\'on: se elimina la escritura defensiva en metodolog\'ia, ceremonias y criterio de aceptaci\'on & Equipo LEIP \\
\end{historial}
\clearpage

% =====================================================================
\section{Introducci\'on}

Este documento tiene como finalidad establecer el plan a utilizar durante el desarrollo del proyecto \proyectoNombre{} -- \proyectoNombreLargo{}. El mismo servir\'a de gu\'ia para establecer las tareas a realizar en cada etapa del proyecto, distribuir responsabilidades entre los integrantes del equipo y definir los hitos que permitir\'an construir y validar la primera etapa del
producto.

El equipo trabajar\'a bajo una metodolog\'ia \'agil (Scrum modificado), con iteraciones cortas (sprints) para la planificaci\'on, desarrollo, validaci\'on y revisi\'on continua del proyecto, manteniendo un tablero de backlog compartido. Las tareas ser\'an divididas y distribuidas entre los integrantes seg\'un las \'areas de responsabilidad de cada uno. Cada integrante sigue el cronograma acordado. Cuando la ejecuci\'on de un sprint obliga a modificarlo, el cambio se documenta y los roles se redistribuyen.

El alcance funcional planificado en este documento corresponde a la primera etapa del producto definida en el Documento de Visi\'on, y comprende: pa\'ises iniciales Brasil, M\'exico y Uruguay; m\'odulos Macro Expectations Intelligence, Prices \& Inflation Intelligence y Real Activity Intelligence; y las capacidades de ingesta, normalizaci\'on, inteligencia anal\'itica, aplicaci\'on web, pagos/suscripciones y API institucional.

% =====================================================================
\section{Gesti\'on del Proyecto}

\subsection{Metodolog\'ia de trabajo}

El proyecto se gestionar\'a mediante un Scrum modificado, adaptado al contexto acad\'emico de la tesis. El trabajo se organiza en sprints de dos semanas, a lo largo de un ciclo total de aproximadamente 8 a 9 meses (entre 16 y 18 sprints), acorde al calendario acad\'emico de la tesis. Al cabo de cada sprint el equipo entrega un incremento funcional del producto y revisa el avance respecto del Product Backlog.

\subsection{Roles del equipo}

Los roles de gesti\'on Scrum se apoyan en la siguiente organizaci\'on de \'areas:

\begin{tabularx}{\textwidth}{@{} >{\bfseries}p{5.4cm} X @{}}
\toprule
\textbf{\'Area / Rol} & \textbf{Responsable(s)} \\
\midrule
Product Owner & Armando Hernandez de la Vega \\
Gesti\'on del proyecto (Scrum Master / coordinaci\'on) & Bruno Olivera, Luana Acu\~na \\
An\'alisis de requerimientos & Armando Hernandez, Bruno Olivera \\
Dise\~no de la interfaz & Bruno Olivera, Sebastian Forische \\
Desarrollo frontend & Luana Acu\~na, Bruno Olivera \\
Desarrollo backend & Sebastian Forische, Luana Acu\~na, Bruno Olivera \\
Infraestructura & Sebastian Forische, Luana Acu\~na \\
QA & Armando Hernandez, Luana Acu\~na \\
Database Administrator & Sebastian Forische \\
Procesamiento de datos con LLM & Armando Hernandez, Sebastian Forische \\
\bottomrule
\end{tabularx}
\bigskip

El resto del equipo participa como equipo de desarrollo multifuncional en las distintas \'areas t\'ecnicas seg\'un la tabla anterior, pudiendo un mismo integrante cubrir varias \'areas seg\'un el sprint.

\subsection{Ceremonias y reuniones}

\begin{itemize}
  \item \textbf{Sprint Planning}: al inicio de cada sprint, para seleccionar y estimar los \'items del backlog a desarrollar.
  \item \textbf{Seguimiento semanal} (daily/weekly stand-up): revisi\'on breve de avance, bloqueos y coordinaci\'on entre \'areas.
  \item \textbf{Sprint Review}: al cierre de cada sprint, demostraci\'on del incremento funcional al equipo y al tutor.
  \item \textbf{Retrospectiva}: identificaci\'on de mejoras al proceso de trabajo y actualizaci\'on del registro de riesgos.
  \item \textbf{Reuni\'on de avance con el tutor de tesis}: semanal. Pasa a quincenal cuando baja la carga de trabajo, acord\'andolo con el tutor.
\end{itemize}

\subsection{Gesti\'on del backlog y priorizaci\'on}

El Product Owner (Armando Hernandez) mantiene y prioriza el Product Backlog junto con el equipo, tomando como insumo los casos de uso especificados y el alcance de la primera etapa del producto. El Product Backlog y las tareas que cada integrante se asigna se gestionan mediante Jira. La priorizaci\'on se basa en: (1) mitigaci\'on de los riesgos m\'as relevantes, (2) dependencias t\'ecnicas entre m\'odulos (ingesta antes que inteligencia, inteligencia antes que visualizaci\'on) y (3) valor para la validaci\'on temprana con usuarios reales.

\subsection{Gesti\'on de cambios}

Todo cambio de alcance, requerimiento o prioridad se documenta como un nuevo \'item o modificaci\'on en el Product Backlog, y se eval\'ua en la siguiente Sprint Planning junto con su impacto en el cronograma. Los cambios que afecten hitos o hip\'otesis clave del proyecto (por ejemplo, restricciones regulatorias sobre una fuente oficial) se reflejan adem\'as como una nueva versi\'on de este Plan de Proyecto, registrada en el Historial de Revisiones.

\subsection{Estimaci\'on}

Los \'items del backlog se estiman en Story Points mediante Planning Poker durante la Sprint Planning. Al finalizar cada sprint se revisa la velocity del equipo, utiliz\'andola como referencia para ajustar el alcance de los sprints siguientes y la proyecci\'on de los hitos definidos en la secci\'on~\ref{sec:hitos}.

\subsection{Gesti\'on de riesgos}

El registro de riesgos se mantiene en la Lista de Riesgos, que es la fuente \'unica de su evaluaci\'on y seguimiento. Este plan no lo duplica: se limita a incorporar su revisi\'on al ritmo de trabajo del equipo.

La Lista de Riesgos se revisa en cada retrospectiva. En esa instancia se verifica si alg\'un riesgo cambi\'o de probabilidad o impacto, si apareci\'o alguno nuevo y si las mitigaciones previstas se implementaron efectivamente. Los riesgos de mayor exposici\'on condicionan directamente la priorizaci\'on del backlog:

\begin{itemize}
  \item \textbf{R-05, condiciones de uso que no habilitan el uso previsto} (nivel cr\'itico): ninguna fuente se incorpora al producto sin su entrada verificada en la matriz de fuentes.
  \item \textbf{R-01, cambios en las fuentes oficiales} (nivel alto): las validaciones autom\'aticas de la ingesta se construyen junto con el primer extractor, no despu\'es.
  \item \textbf{R-03, errores o inconsistencias en los datos} (nivel alto): la normalizaci\'on y la cuarentena preceden a cualquier trabajo sobre la capa de inteligencia.
  \item \textbf{R-09, baja adopci\'on del p\'ublico objetivo} (nivel alto): la validaci\'on con usuarios reales se sostiene a lo largo de todo el ciclo y no se concentra al final.
\end{itemize}

% =====================================================================
\section{Objetivos e hitos del proyecto}
\label{sec:hitos}

A continuaci\'on se definen los objetivos a alto nivel de cada hito del proyecto, agrupando los sprints del Scrum modificado del equipo seg\'un los bloques funcionales de la primera etapa del producto. Cada hito cierra con una versi\'on funcional (release) del incremento correspondiente, demostrada al equipo y al tutor de tesis.

\renewcommand{\arraystretch}{1.2}
\begin{longtable}{@{}>{\raggedright\arraybackslash}p{2.6cm} >{\raggedright\arraybackslash}p{2cm} >{\raggedright\arraybackslash}p{\dimexpr\textwidth-2.6cm-2cm-6\tabcolsep\relax}@{}}
\toprule
\textbf{Hito} & \textbf{Sprints} & \textbf{Objetivos y entregables principales} \\
\midrule
\endhead
\bottomrule
\endfoot

\textbf{Hito 0}\newline Fundacional & Sprint 0 &
\begin{itemize}[leftmargin=1em]
  \item Definici\'on y validaci\'on de la arquitectura general (Node.js, Python, React, PostgreSQL, Azure Blob Storage, Azure Container Apps).
  \item Configuraci\'on de repositorios en GitHub, pipeline de CI/CD con GitHub Actions y entornos de desarrollo.
  \item Definici\'on del modelo de datos inicial y del patr\'on Medallion (Bronze/Silver/Gold/Cuarentena).
  \item Elaboraci\'on y priorizaci\'on del Product Backlog inicial en base a los casos de uso identificados.
  \item Mitigaci\'on temprana de los riesgos t\'ecnicos y legales identificados como prioritarios (uso de fuentes oficiales, formatos heterog\'eneos).
  \item \textbf{Hito}: acuerdo formal del alcance de la primera etapa, validado con el tutor de tesis.
\end{itemize} \\
\midrule

\textbf{Hito 1}\newline Ingesta y normalizaci\'on & Sprints 1 a 4 &
\begin{itemize}[leftmargin=1em]
  \item Desarrollo de los pipelines de ingesta de fuentes oficiales de Brasil (Boletin Focus -- BCB), M\'exico (Encuesta de expectativas / SIE -- Banxico) y Uruguay (Encuestas de expectativas del BCU).
  \item Implementaci\'on del extractor estructurado basado en LLM para normalizar HTML/archivos heterog\'eneos a JSON interno.
  \item Registro de fuente, fecha de publicaci\'on, fecha de ingesta, formato, versi\'on y checksum de cada dato (capa Bronze).
  \item Normalizaci\'on de variables, unidades, fechas y frecuencias hacia el modelo com\'un (capa Silver).
  \item Control de calidad y trazabilidad: detecci\'on de duplicados/faltantes/inconsistencias y env\'io a cuarentena.
  \item \textbf{Meta}: relevar y detallar al menos el 80\% de los casos de uso identificados.
  \item \textbf{Hito}: al menos 1 caso de uso representativo (ingesta y normalizaci\'on de una fuente) funcionando end-to-end, demostrado al tutor.
\end{itemize} \\
\midrule

\textbf{Hito 2}\newline Capa de inteligencia & Sprints 5 a 7 &
\begin{itemize}[leftmargin=1em]
  \item Motor determin\'istico de m\'etricas: cambios, momentum, divergencia regional, percentiles hist\'oricos, sorpresa, persistencia y distancia frente a metas.
  \item Grafo de dependencias para cruces entre m\'odulos (inflaci\'on esperada vs. observada; crecimiento esperado vs. actividad efectiva).
  \item Generaci\'on autom\'atica del resumen descriptivo del radar mediante LLM a partir de m\'etricas ya calculadas (capa Gold).
  \item Persistencia de m\'etricas, cruces y se\~nales en PostgreSQL con trazabilidad completa.
  \item \textbf{Hito}: al menos 3 casos de uso representativos (ingesta, c\'alculo de m\'etricas y cruce entre m\'odulos) funcionando end-to-end, demostrados al tutor.
\end{itemize} \\
\midrule

\textbf{Hito 3}\newline Aplicaci\'on web core & Sprints 8 a 11 &
\begin{itemize}[leftmargin=1em]
  \item Autenticaci\'on, roles, permisos y entitlements por plan, m\'odulo, pa\'is y funcionalidad.
  \item Dashboard de an\'alisis: selecci\'on de pa\'is/m\'odulo/variable/fuente/per\'iodo, gr\'aficos, tablas, rankings e hist\'oricos.
  \item Radar de expectativas con versi\'on resumida para el acceso sin suscripci\'on y versi\'on completa para el acceso contratado, incluyendo alertas y briefings.
  \item Exportaci\'on de resultados en Excel/PDF y guardado de consultas frecuentes.
\end{itemize} \\
\midrule

\textbf{Hito 4}\newline Monetizaci\'on e integraciones & Sprints 12 a 14 &
\begin{itemize}[leftmargin=1em]
  \item Integraci\'on de pagos y suscripciones mediante checkout externo y webhooks (sin almacenar datos de tarjetas).
  \item API institucional autenticada, con cuotas y permisos seg\'un plan contratado.
  \item Aislamiento multi-tenant para fuentes privadas de las organizaciones cliente.
  \item Panel administrativo para gesti\'on de usuarios, organizaciones, planes y suscripciones.
\end{itemize} \\
\midrule

\textbf{Hito 5}\newline Validaci\'on, QA y cierre & Sprints 15 y 16 &
\begin{itemize}[leftmargin=1em]
  \item Pruebas end-to-end de los flujos cr\'iticos (ingesta $\rightarrow$ normalizaci\'on $\rightarrow$ inteligencia $\rightarrow$ visualizaci\'on $\rightarrow$ exportaci\'on).
  \item Validaci\'on del producto con usuarios reales del segmento objetivo (continuaci\'on del plan de validaci\'on iniciado con la encuesta a analistas).
  \item Ajustes finales de producto, seguridad y performance seg\'un resultados de QA y validaci\'on.
  \item Actualizaci\'on de la documentaci\'on con el estado final del producto y preparaci\'on de la defensa de tesis.
\end{itemize} \\

\end{longtable}
\renewcommand{\arraystretch}{1}

% =====================================================================
\section{Criterio de Aceptaci\'on del Proyecto}

El criterio de aceptaci\'on del proyecto se define en base al alcance de la primera etapa del producto y al plan de validaci\'on. Este criterio se valida con el tutor de tesis y con las organizaciones que participen de la validaci\'on del producto.

\textbf{El proyecto se acepta cuando se cumplen las siguientes condiciones:}

\begin{enumerate}
  \item Ingesta y normalizaci\'on funcionando de forma automatizada para las fuentes oficiales de Brasil, M\'exico y Uruguay, para los m\'odulos Macro Expectations Intelligence, Prices \& Inflation Intelligence y Real Activity Intelligence.
  \item C\'alculo autom\'atico de las m\'etricas determin\'isticas definidas (cambios, momentum, divergencia regional, percentiles hist\'oricos, sorpresa, persistencia, distancia frente a metas) y de los cruces iniciales entre m\'odulos.
  \item Aplicaci\'on web operativa con autenticaci\'on, roles/entitlements, dashboard de an\'alisis, radar de expectativas (versi\'on resumida y versi\'on completa) y exportaci\'on de resultados.
  \item Trazabilidad end-to-end: cada dato, m\'etrica y exportaci\'on conserva fuente, fecha, versi\'on y criterio de procesamiento.
  \item Flujo de pagos/suscripciones y API institucional funcionando de forma segura (sin almacenamiento directo de datos de tarjetas), con al menos un caso de uso institucional validado.
  \item Validaci\'on del problema y del producto con usuarios reales del segmento objetivo, dando continuidad a la encuesta inicial ya realizada, con resultados documentados.
\end{enumerate}

\end{document}
